\section{Experiments and Discussions}
\label{section:exp}
\input{tabs/relatedWorksComparison}

\subsection{Experimental Settings}

\noindent\textbf{Benchmark.}
We evaluate on PIE-Bench~\cite{PIEBench}, a comprehensive benchmark comprising 700 images (both natural and artificial) across 10 editing categories, including object replacement, addition, removal, attribute modification, and style transfer.

\noindent\textbf{Metrics.}
Following prior work~\cite{PIEBench,AREdit}, we adopt six metrics grouped into two aspects. For \textit{background preservation}, we report Structure Distance~\cite{PIEBench} (S.D.), PSNR, SSIM, and LPIPS~\cite{LPIPS}. For \textit{editing quality}, we report CLIP similarity~\cite{CLIP} on both the whole image (CLIP whole) and the edited region (CLIP edited). All metrics are computed using the official PIE-Bench evaluation pipeline. Inference time per image is also reported.

\noindent\textbf{Baselines.}
We compare with 11 training-free editing methods spanning three architectural families: six diffusion-based — Prompt-to-Prompt~\cite{prompt2prompt}, Pix2Pix-Zero~\cite{pix2pix-zero}, MasaCtrl~\cite{masactrl}, PnP~\cite{PnP}, PnP Inversion~\cite{PnP_Inversion}, and LEDits++~\cite{LEDits++} — four flow-based — FlowEdit~\cite{Flowedit}, RF-Inversion~\cite{RF-inversion}, ReFlex~\cite{ReFlex}, and FireFlow~\cite{FireFlow} — and one autoregressive method, AREdit~\cite{AREdit}. Baseline results are obtained using officially released code under default settings, except AREdit whose results are cited directly from the original paper as no official implementation has been released.

\noindent\textbf{Implementation details.}
We build upon Infinity~\cite{Infinity} (2B parameters) with BSQ-VAE~\cite{BSQVAE} ($d{=}32$). Generation uses $S{=}13$ scales ($1{\times}1$ to $64{\times}64$), CFG scale 4.0, and temperature $\tau{=}0.5$. Text conditioning uses Flan-T5-XL~\cite{flan-t5}. We set $N{=}2$ full-replace scales. For binarization, we use adaptive thresholding that adjusts the percentile parameter $p$ based on the attention distribution of each image (see Sec.~\ref{sec:mask_construction}). For IQR filtering, we discard transformer blocks whose attention deviation from the cross-block mean exceeds the IQR fence (Sec.~\ref{sec:mask_construction}), using blocks 2--31 of the 32-block model (blocks 0--1 are excluded as they consistently exhibit the global attention effect described in Sec.~\ref{sec:mask_construction}). All experiments run on a single NVIDIA RTX 5090 GPU with ${\sim}$2.5\,s inference per image.


\subsection{Comparison with State-of-the-Art Methods}

\noindent\textbf{Quantitative results.}
Table~\ref{tab:method_comparison_wo_kvedit} summarizes the comparison. Our method achieves the best SSIM (0.8565) and LPIPS (0.0833) among all methods, and the second-best CLIP whole (0.2613) and CLIP edited (0.2290), demonstrating strong background preservation and competitive edit quality with only 2B parameters — 6$\times$ fewer than all flow-based baselines.

\textit{vs.\ diffusion-based methods.}
Compared to the best diffusion-based method PnP Inversion, we improve SSIM by 7.5\% and LPIPS by 21.5\% while achieving higher CLIP scores on both whole-image and edited-region metrics.

\textit{vs.\ flow-based methods.}
Against the strongest flow-based competitor ReFlex, we attain better SSIM and LPIPS despite the 6$\times$ parameter gap, though ReFlex achieves higher PSNR.

\textit{vs.\ AREdit.}
Compared to AREdit~\cite{AREdit}, a VAR-based method with the same 2B backbone, we achieve uniformly better metrics across all six measures, validating the effectiveness of attention-driven mask construction over cached token statistics.

We report inference time for reference, but note that methods operate at different resolutions.

\input{imgs/qualitative}

\noindent\textbf{Qualitative results.}
Fig.~\ref{fig:qualitative} presents representative visual comparisons. Diffusion-based methods often suffer from structural drift or incomplete edits due to inversion errors: P2P and Pix2Pix-Zero alter the global layout, while MasaCtrl preserves structure but produces blurry edited regions. Among flow-based methods, ReFlex achieves the best preservation but occasionally misses fine-grained edit details. In contrast, our method produces clean, precisely localized edits while faithfully preserving the background, benefiting from explicit attention-driven masking and discrete token replacement.


\subsection{Ablation Studies}
\label{sec:ablation}

\noindent\textbf{Effect of attention-based masking.}
We compare our full pipeline against a variant that disables attention masks (Scale-$N$): the pipeline is identical to ours, but only the first $N$ scales share source tokens via full replacement while all subsequent scales generate freely under $P_t$ with no spatial guidance. We vary $N \in \{1, 2, 4, 6\}$.

\input{tabs/p2pAndAttnMask}

As shown in Table~\ref{tab:ablation_mask}, increasing shared scales progressively improves preservation, but gains saturate: even at $N{=}6$ (nearly half of 13 scales), SSIM reaches only 0.553, because scales beyond $N$ generate independently with no source reference. Our attention-masked method provides per-token spatial guidance at every scale, improving SSIM by 55\% and LPIPS by 67.5\% over the best Scale-$N$ variant, demonstrating that spatially precise masks are far more effective than indiscriminate scale-level replacement.
Qualitatively, few shared scales degrade the background, while many shared scales prevent editing since the global layout is already locked. VAR commits to structure in the earliest scales, so shape-changing edits require modifying early tokens — which uniform replacement cannot accommodate. Our attention-mask approach resolves this tension by selectively replacing only background tokens, achieving both faithful editing and preservation. We provide detailed visual analysis in the supplementary material.

%  ========== phase 2 ablation
\noindent\textbf{Effect of Phase~2.}
Phase~2 performs a second generation pass under the target prompt to extract the target focus mask $M_{k,t}^{\mathrm{f}}$ before token replacement (Sec.~\ref{sec:pipeline}). We compare our full three-phase pipeline against a variant that skips Phase~2.

\input{tabs/ablationPhase2}

As shown in Table~\ref{tab:ablation_phase2}, Phase~2 yields the largest single-component improvement: PSNR increases by 37.7\% and SSIM by 37.5\%, while LPIPS drops by 59.3\% and Structure Distance decreases by 39.0\%. CLIP scores also improve by 3.5--3.6\%. These results confirm that the source-side mask alone cannot accurately delineate editing boundaries — particularly when the target introduces new spatial content absent in the source. The dual-stream mask construction of Phases~1--2 is essential for preventing background drift while enabling precise editing.

\noindent\textbf{Effect of adaptive thresholding.}
A key hyperparameter in our pipeline is the binarization threshold that converts continuous attention maps into discrete edit/preserve masks. We compare our adaptive thresholding strategy — which adjusts the percentile parameter based on the attention distribution of each image — against fixed-percentile thresholds at 70\%, 80\%, and 90\%.

\begin{table}[t]
\centering
\caption{Ablation studies on attention mask construction (AMC). Best in \textbf{bold}.}
\label{tab:ablation_amc}
\small
\resizebox{\linewidth}{!}{%
\begin{tabular}{l|cccc|cc}
\toprule
Method & S.D.$\downarrow$ & PSNR$\uparrow$ & SSIM$\uparrow$ & LPIPS$\downarrow$ & CLIPw$\uparrow$ & CLIPe$\uparrow$ \\
\midrule
Adaptive          & \textbf{0.0316} & \textbf{23.44} & \textbf{0.8565} & \textbf{0.0833} & \textbf{0.2613} & \textbf{0.2290} \\
90\% (fixed)     & 0.0484 & 23.01 & 0.8357 & 0.0945 & 0.2578 & 0.2246 \\
80\% (fixed)     & 0.0482 & 22.15 & 0.8029 & 0.1024 & 0.2537 & 0.2205 \\
70\% (fixed)     & 0.0489 & 21.02 & 0.7684 & 0.1135 & 0.2468 & 0.2157 \\
\bottomrule
\end{tabular}%
}

\vspace{2pt}
\centerline{(a) Adaptive vs.\ fixed-percentile thresholding.}

\vspace{6pt}

\resizebox{\linewidth}{!}{%
\begin{tabular}{l|cccc|cc}
\toprule
Method & S.D.$\downarrow$ & PSNR$\uparrow$ & SSIM$\uparrow$ & LPIPS$\downarrow$ & CLIPw$\uparrow$ & CLIPe$\uparrow$ \\
\midrule
With IQR  & \textbf{0.0316} & \textbf{23.44} & \textbf{0.8565} & \textbf{0.0833} & \textbf{0.2613} & \textbf{0.2290} \\
W/O IQR   & 0.0328 & 20.56 & 0.7872 & 0.1285 & 0.2550 & 0.2260 \\
\bottomrule
\end{tabular}%
}

\vspace{2pt}
\centerline{(b) Ablation on IQR filtering.}

\end{table}


As shown in Table~\ref{tab:ablation_amc}(a), adaptive thresholding outperforms all fixed-percentile variants across every metric, achieving the lowest Structure Distance (0.0316 vs.\ 0.0482--0.0489) and the best preservation and edit quality scores. Among fixed variants, performance degrades monotonically from 90\% to 70\%, confirming that no single percentile universally accommodates the diversity of attention distributions across images and edit types.

\input{imgs/ablationDynamic}

Fig.~\ref{fig:ablationDynamic} illustrates this qualitatively: adaptive thresholding produces fine-grained masks that adapt to each image's attention distribution, preserving subtle details such as thin branches. Fixed thresholds at 70--80\% yield overly coarse masks with significant background structure loss, while 90\% is reasonable but still misses fine details.

\noindent\textbf{Effect of IQR filtering.}
Not all transformer blocks produce equally informative attention maps. As described in Sec.~\ref{sec:mask_construction}, we measure each block's deviation from the cross-block mean and discard outlier blocks whose deviation exceeds the IQR fence. We compare our full model (IQR-filtered, blocks 2--31) against naively averaging all 32 blocks.

As shown in Table~\ref{tab:ablation_amc}(b), IQR filtering yields substantial improvements: PSNR increases by 14.0\% and SSIM by 8.8\%, while LPIPS drops by 35.2\%, confirming that filtered masks avoid corrupting non-target regions. Even Structure Distance improves (0.0328$\to$0.0316), showing that cleaner masks also reduce unintended structural changes.

% Content moved to ablationPhase2.tex and ablationThreshold.tex

Fig.~\ref{fig:ablationIQR} illustrates this: without filtering, noisy attention from uninformative blocks produces imprecise masks that distort the bicycle structure; with IQR filtering, the mask cleanly delineates the target, preserving fine structural details.

